\emph{(Preliminaries: the two facts of the field.)} In the Painlev\'e-Gullstrand gauge the physical coframe is $\theta^{0} = dt$, $\theta^{1} = \beta\,dt + dr$ (Proposition \claimref{singularity.connection-bivector-computed}), so the background direction $e_{t} + w e_{r}$ has physical image $\gamma_{0} + (\beta + w)\gamma_{1}$, null for $w = 1 - \beta$ (outgoing) and $w = -(1 + \beta)$ (ingoing). For the ingoing tangent $a = e_{t}/(1 + \beta) - e_{r}$ the image is $k = (\gamma_{0} - \gamma_{1})/(1 + \beta)$ and $dr/d\lambda = -1$. Then $a\cdot\nabla k = \beta'(1 + \beta)^{-2}(\gamma_{0} - \gamma_{1})$, $\Omega(a) = (1+\beta)^{-1}\Omega(e_{t}) - \Omega(e_{r}) = \beta'(1 + \beta)^{-1}\sigma_{r}$, and with $\sigma_{r} = \gamma_{1}\gamma_{0}$ one has $\sigma_{r}\times(\gamma_{0} - \gamma_{1}) = -(\gamma_{0} - \gamma_{1})$, so $a\cdot\nabla k + \Omega(a)\times k = 0$: $k$ is parallel and $r$ is affine (Fact \claimref{synthesis.imported-gtg-dictionary}(iii)). The background time translation is a symmetry of the field ($\beta$ depends on $r$ only) and moves the outgoing worldlines; label them by $u = t - r_{*}$ with $dr_{*}/dr = 1/(1 - \beta)$. Along an ingoing worldline $dt/dr = -1/(1 + \beta)$, so $du/dr = -2/(1 - \beta^{2}) = -2r/(r - 2M)$. For $U = -\zeta e^{\zeta/2M}$, $dU/du = -(1 + \zeta/2M)e^{\zeta/2M}\,dr/du = \tfrac{2M + \zeta}{2M}e^{\zeta/2M}\cdot\tfrac{\zeta}{2(2M + \zeta)} = \tfrac{1}{4M}\zeta e^{\zeta/2M} = -\kappa U$. Hence $U(u + t) = e^{-\kappa t}U(u)$ on each ingoing ray, and $dU/dr = -(1 + \zeta/2M)e^{\zeta/2M} = -1$ at $\zeta = 0$. The check verifies the null directions, the parallel transport (with a control in which $\Omega$ is reversed), $dr/du$ and $dU/du = -\kappa U$ exactly at $(M, \beta)$ and $(M, r)$ samples, horizon and interior included.

\emph{(i)} Throughout, the modes are the per-light-ray chiral fields of the s-wave reduction of the postulate J3 (\claimref{hawking.joint-j3-s-wave-reduction}) in its geometric-optics form; this is a hypothesis of the result. Inside the shell the field vanishes: an ingoing ray of advanced time $v$ reaches the centre and leaves as the outgoing ray of flat retarded time $v$, which meets the shell $v = v_{0}$ at $r = s = (v_{0} - v)/2$. The radial label is continuous across the shell, so in the Schwarzschild region the ray is the outgoing ray through $(v_{0}, s)$; in the derived-metric reading its retarded time is $u = v_{0} - 2r_{*}(s)$, $r_{*}(s) = s + m\ln(s/m - 1)$ (Fact \claimref{synthesis.imported-gtg-dictionary}), and its Kruskal label is $-4M e^{-v_{0}/4M}e^{s/m}(s/m - 1) = C\,z e^{z}$ with $z = s/m - 1$, $C = -4Me^{1 - v_{0}/4M}$; the escaping rays are $s > m$, the horizon ray $s = m$ ($v = v_{H}$), and the formula continues analytically to $s < m$. On the outgoing modes of an ingoing ray crossing the future horizon the label $U$ is the analytic coordinate of (i), and by the preliminaries $U$ is an analytic function of $r$ with $dU/dr = -1$ there; so the parameter $v$ of the in-vacuum is an analytic function $f$ of the affine parameter $\lambda$ of the ray with $f' \ne 0$ at the crossing. The in-vacuum's kernel in $v$ is the null vacuum (hypothesis), and its pull-back to $\lambda$ is $f'(\lambda_{1})^{1/2}f'(\lambda_{2})^{1/2}/(f(\lambda_{1}) - f(\lambda_{2}))$. For any analytic $f$ with $f'(0) \ne 0$ this equals $(\lambda_{1} - \lambda_{2})^{-1} + \tfrac{1}{12}\{f, \lambda\}(\lambda_{1} - \lambda_{2}) + \dots$, the remainder analytic and zero on the diagonal (the weight-$\tfrac12$ kernel is exactly M\"obius invariant, so only the Schwarzian enters at the first order). The check computes the series exactly for $f(\lambda) = \lambda e^{\lambda}$, the common form of the ray's Kruskal map and the ray-tracing map after affine changes (which the kernel ignores up to the M\"obius invariance it verifies), along the lines $(\lambda_{1}, \lambda_{2}) = (p\epsilon, q\epsilon)$ for three $(p, q)$, to order six: the $\epsilon^{0}$ term of $\epsilon R$ vanishes, the $\epsilon^{1}$ term vanishes, and the $\epsilon^{2}$ term is $-\tfrac{3}{12}(p - q)$, with $\{\lambda e^{\lambda}, \lambda\}_{0} = -3$; a control with $\lambda e^{2\lambda}$ gives another coefficient. Composition of analytic maps with nonzero derivative keeps the class, so the in-vacuum is in the short-distance class at the crossing. In closed form $\{U, y\} = (y + 2)/y - \tfrac32((y + 1)/y)^{2}$ with $y = s/m$, singular only at $s = 0$, the ray that reaches the shell at the centre, which is inside the horizon; at the horizon ray $\{U, v\} = \{U, y\}/(4m^{2}) = -3/(4m^{2})$ (checked).

\emph{(ii)} For a chiral field the vacuum-subtracted null energy in the parameter $u$ of a state that is the vacuum in the parameter $v$ is $-\tfrac{c}{24\pi}\{v, u\}$, with $c = 1$ for a complex chiral fermion; this is the coefficient of the linear remainder in (ii) under point splitting (cited from memory: Davies and Fulling 1977, the moving mirror; the same law for every two-dimensional conformal field). With $v = v_{0} - 2s$ and $u'(s) = -2s/(s - m)$, $\{v, u\} = -\{u, s\}/u'(s)^{2}$, and $\{u, s\} = m(4s - 3m)/(2s^{2}(s - m)^{2})$, so $\{v, u\} = -m(4s - 3m)/(8s^{4})$. The check obtains $u''$ and $u'''$ as exact Taylor coefficients of the rational function $u'$ at twelve sample $(m, s)$ and compares. At $s = m$ the closed form is $-1/(8m^{2}) = -\kappa^{2}/2$; as $u \to \infty$, $s/m - 1 \sim C'e^{-\kappa u}$, and the deviation is $O(s - m)$ (checked at $s - m = m/10, m/1000$). So the late flux is $\kappa^{2}/48\pi$. A complex chiral fermion at temperature $T$ carries $2\int_{0}^{\infty}\varpi(1 + e^{\varpi/T})^{-1}d\varpi/2\pi = \pi T^{2}/12$ (the Fermi integral $\pi^{2}/12$, cited); at $T = T_{H}$ this is $\kappa^{2}/48\pi$. For the two-point function: with $x = e^{\kappa u/2}$, $U = -1/(\kappa x^{2})$, $U' = x^{-2}$, and $(U_{1}'U_{2}')^{1/2}/(U_{1} - U_{2}) = \kappa/(x_{1}/x_{2} - x_{2}/x_{1}) = \tfrac{\kappa}{2}\sinh(\tfrac{\kappa}{2}(u_{1} - u_{2}))^{-1}$ exactly; continuation by $-I_{4}\beta$, $\beta = 2\pi/\kappa$ (numbers, the kernels' values included, are read in $\mathrm{span}\{1, I_{4}\}$, $I_{4}$ the field theory's unit), sends $x_{1} \mapsto -x_{1}$ and the kernel to minus itself, with no pole on the strip, the fermionic KMS boundary value. By Theorem \claimref{hawking.i-state-continuum} this quasi-free state is the continuum I-state for the Killing translation at $\beta$, and it is the unique $\beta$-KMS state by Fact \claimref{hawking.imported-ideal-fermi-gas-uniqueness}. Equivalently, the half-line algebra $U < 0$ of the null vacuum carries the translations $U \mapsto U + a$, $a \ge 0$, as a half-sided inclusion with positive generator, and the abstract theorem of Fact \claimref{angle.imported-theorems-deliver-their-hypotheses} makes its modular group the dilations, which by (i) are the Killing translations at rate $\kappa$. The remainder: at Killing shift $t$ the kernel in $u$ is $e^{-\kappa t}$ times the $U$-kernel at points scaled by $\mu = e^{-\kappa t}$; the vacuum part is exactly scale invariant and the remainder contributes $\mu R(\mu\lambda_{1}, \mu\lambda_{2}) = O(\mu^{2})$ (checked at $\mu = 1/10, 1/100$).

\emph{(iii)} The transformation law gives $F_{u} = U_{u}^{2}F_{U} - \tfrac{1}{24\pi}\{U, u\}$, and for $U = -e^{-\kappa u}/\kappa$, $\{U, u\} = -\kappa^{2}/2$ ($(\log U')' = -\kappa$), $U_{u} = -\kappa U$. A state of the outgoing modes that is KMS for the Killing translations at late times at temperature $\theta T_{H}$ has late flux $\tfrac{\pi}{12}\theta^{2}T_{H}^{2}$ per complex chiral field (the Fermi integral above). A late flux $\tfrac{\pi}{12}\theta^{2}T_{H}^{2} = \theta^{2}\kappa^{2}/48\pi$ therefore needs $F_{U} = (\theta^{2} - 1)/(48\pi U^{2})$ to leading order, bounded as $U \to 0^{-}$ only at $\theta = 1$; in Schwarzian units ($F = -S/24\pi$) the check verifies the closure of the chain rule and the dichotomy at four $\theta$, two $\kappa$, three $U$. The Boulware state, the vacuum for $u$, has $u = -\kappa^{-1}\ln(-\kappa U)$, $(\log|u'|)' = -1/U$, and $\{u, U\} = 1/(2U^{2})$: its energy in the horizon-regular frame diverges at the future horizon. This is the Candelas / Christensen-Fulling regularity criterion on the series' carrier. The same formula, with $\theta = 1/c$, is the flux of a state that is stationary for a boost at period $2\pi c$ near the edge of a horizon generator (filed alongside, synthesis.stationary-horizon-state-is-its-own-edge-scaling-limit); the late value $\kappa^{2}/48\pi$ is the Robinson-Wilczek anomaly flux (cited from memory).

\emph{(Remark, from the result filed alongside, synthesis.collapsing-hole-flux-from-ground-state-on-the-ingoing-sheet.)} In the collapse field the future horizon is, inside the shell, the outgoing light cone of a point of the centre, so its generators begin there; an ingoing radial ray ends at $r = 0$ at finite affine parameter. On $L^{2}((0, \infty))$ the solutions of $I_{4} f' = \pm I_{4} f$ are $e^{\lambda}$ and $e^{-\lambda}$, of which only $e^{-\lambda}$ is square integrable: deficiency indices $(1, 0)$, so $I_{4}\,d/d\lambda$ has no self-adjoint extension and generates no unitary translation group (von Neumann; cited). An exact ground-state condition of the form of (N$_{+}$) has nothing to be a ground state for, and only the short-distance form survives. Combining (i), (ii) and (iii): the collapse realises the short-distance class, and among states KMS for the Killing translations at late times finite null energy across the future horizon holds only at $T = T_{H}$; bounded $F_{U}$ is necessary for the class and not sufficient, so (iii) does not characterise the class. What is not shown: the four-dimensional statement with greybody factors and angular modes (Fredenhagen and Haag 1990, cited from memory, not imported), and other collapse geometries (argued: any collapse in which the ray tracing is $C^{3}$ with nonzero derivative at the horizon ray gives the same limit, since only $\{U, v\}$ bounded is used).
